% === IX. CONCLUSION ==========================================================
\section{Conclusion and Future Work}
\label{sec:conclusion}

This paper studied the identifiability of a learning-based on-track Pacejka
identification pipeline~\cite{ontrack2025} at full vehicle scale, in a CARLA
architecture exposing the plant's per-wheel forces as ground truth.
Transplanted unchanged from a 1:10 platform it fails in a specific, diagnosable
and correctable way, and not through the solver, the bounds or the tuning. The
virtual-data step's own reachable friction
$\mu_{\mathrm{reach}}=v^{2}\delta_{\mathrm{sweep}}/(gL)$ depends on the rollout
configuration and not on the tire, and a rollout speed appropriate at 1:10 caps
it at \num{0.43} at full scale, below the plant's peak, so the peak factor is
not
identifiable from the synthetic data and the bounded optimiser resolves the
ambiguity on a box edge.

The correction has six parts: an excitation-aware rollout speed, an
extrapolation bound on the sweep, a frozen regularisation target,
identification against the achieved road-wheel angle, measured mass and
geometry, and a friction warm start read from the curvature of the axle force
curve and applied only as a floor. Together they produce identifications with
no railed coefficient, a peak factor of \num{1.009}/\num{1.002} against a plant
realising \num{1.05}, and the correct understeer sign. Over three timed laps on
a static surface the pipeline \emph{preserves} that peak factor to
\num{0.041}/\num{0.048} RMSE where the inherited configuration starts and ends
on its \num{2.0} box constraint, and halves the MPC's lateral tracking error,
\SI{0.177}{\meter} to \SI{0.092}{\meter} RMS, at unchanged solver cost.
Preserves, not identifies: this trajectory never approaches the tire's peak and
the accepted fits return the configured prior to four decimals. Recovery
\emph{when the excitation reaches the peak} is shown on a known tire in
Section~\ref{subsec:sweep}. Each configuration is a single run, so the
closed-loop margins are observations rather than estimates with a confidence
interval. We further argue that per-coefficient box constraints are not a
sufficient admissibility criterion for handing a model to a model-based
controller, and gate on derived physical quantities at that interface
instead.

Equation~\eqref{eq:mureach} inverts into a requirement on the manoeuvre, so a
controller that inserts bounded slalom excitation when the identifier reports
insufficient slip coverage would move the pipeline toward active,
safety-bounded input design~\cite{goodwin1977design}. The decisive test of
every claim here is a physical vehicle: the identification code is unchanged
between simulation and hardware, and only the ground truth goes away.
